\section{Depth guarantees and their limits}\label{dp:section}

An orbit set must accommodate all directions of the corner with one convex
quadratic geometry. The relevant scale is the size of the corner at its
optimal objective level relative to the violation at its vertex. This scale
gives a uniform approximation guarantee. It also identifies a regime in which
even completed orbit sets can lose almost the entire corner bound.

Throughout this section, $S=\{(x,y,w):q(x,y,w)=w-xy\le0\}$,
$q_0=q(\sbar)>0$, there are finitely many rays $p_j$, and
$\omega_j>0$. We assume that the feasible corner is nonempty, so that
$0<\zK<\infty$. Define
\begin{equation}\label{dp:scaled-rays}
 \widetilde p_j=\frac{\zK}{\omega_j}p_j,\qquad
 T_r=\conv\{\sbar,\sbar+r\widetilde p_j:1\le j\le N\},\qquad T^*=T_1.
\end{equation}
Every point of $T_r$ has a representation of objective value at most
$r\zK$. Hence $q>0$ on $T_r$ when $r<1$, and $q\ge0$ on $T^*$.
An admissible convex set has cut ratio at least $r$ precisely when its
steps along all scaled rays are at least $r$, equivalently when it contains
$T_r$. We use $\zA$ and $\zB$ for the suprema over the orbit sets and their
closed upward completions, respectively, with $\sbar$ in their interiors.

\subsection{The scale of the optimal simplex}

Put
\begin{equation}\label{dp:depth}
 \widetilde X=\max_j|\widetilde p_{jx}|,\qquad
 \widetilde Y=\max_j|\widetilde p_{jy}|,\qquad
 D=\sqrt{\frac{\widetilde X\widetilde Y}{q_0}}.
\end{equation}
Thus $D^2$ compares the largest possible bilinear curvature scale of the
scaled rays with the violation at the vertex. Its reciprocal
$\delta=1/D^2$ is a relative depth, with $\delta=\infty$ when $D=0$.

\begin{lemma}[Symmetries of depth]\label{dp:symmetries}
For $ab>0$, the affine maps
\[
 (x,y,w)\longmapsto
 (ax+u,by+v,ab\,w+avx+buy+uv)
\]
and the interchange of $x$ and $y$ preserve $S$, transport both families
onto themselves, and preserve $\zK,\zA,\zB$ when the rays are transported
and the costs are kept fixed. They preserve $D$. The relative discriminant
\begin{equation}\label{dp:relative-discriminant}
 \Delta_{\mathrm{rel}}=
 \frac{B^2-4Aq_0}{B^2+|4Aq_0|}
\end{equation}
of a restriction $q(\sbar+sp)=As^2+Bs+q_0$ is also preserved whenever its
denominator is nonzero. Euclidean angles, the ray condition number, and the
fixed rule studied below need not be preserved.
\end{lemma}

\begin{proof}
The map multiplies $q$ by $ab>0$. In matrix coordinates it is
$M'=L M R^T$, with
$L=\left(\begin{smallmatrix}a&u\\0&1\end{smallmatrix}\right)$ and
$R=\left(\begin{smallmatrix}b&v\\0&1\end{smallmatrix}\right)$.
Writing $G=F^T$, the transformed orbit matrix is $G'=RGL^{-1}$, since
$\sym(G'M')=R\sym(GM)R^T$. Its determinant is positive. The map sends
$e_w$ to $ab e_w$, so it also transports upward completions. Transposition
uses $G'=G^{-1}$ and a congruence by $G$. Ray coordinates and their objective
values are unchanged. Finally $\widetilde X,\widetilde Y,q_0$ are multiplied
by $|a|,|b|,ab$, respectively; the coefficients in the ray restriction share
the multiplier $ab$.
\end{proof}

These are symmetries of the bilinear model, rather than arbitrary affine
changes of coordinates. A constant ray restriction has no defined relative
discriminant. A zero discriminant means a repeated root; a grazing contact
on the forward ray additionally requires $B<0$. These qualifications avoid
using a ray angle as a substitute for vertex depth.

For $t>0$, the parabolic cylinder centered at the vertex is
\begin{equation}\label{dp:cylinder}
 C_t=\left\{s:q(s)\ge
 \frac{\bigl(t(x-\bar x)-(y-\bar y)/t\bigr)^2}{4}\right\}.
\end{equation}
It is the epigraph of a convex quadratic and the orbit set with
$G=F^T=\left(\begin{smallmatrix}t&-t\bar x+\bar y/t\\0&1/t\end{smallmatrix}\right)$.
Indeed its determinant inequality is \cref{dp:cylinder}, and the second
diagonal entry of its PSD matrix is the positive constant $1/t$.
It equals its own upward completion, and its interior contains $\sbar$.
Along a scaled ray, define
\[
 \ell_j=\frac{\nabla q(\sbar)^T\widetilde p_j}{q_0},\qquad
 b_j(t)=\frac{t\widetilde p_{jx}+\widetilde p_{jy}/t}{\sqrt{q_0}}.
\]
The cylinder slack, divided by $q_0$, is
$1+\ell_js-b_j(t)^2s^2/4$, and its step is
\begin{equation}\label{dp:cylinder-step}
 \alpha_j(t)=\frac{2}{\sqrt{\ell_j^2+b_j(t)^2}-\ell_j}.
\end{equation}
A zero denominator means an infinite step. Thus a single scalar parameter
$t$ already gives a cut with a useful uniform bound.

\begin{theorem}[Depth guarantee]\label{dp:depth-bound}
With the hypotheses above, without a rank assumption on the projected rays,
\begin{equation}\label{dp:depth-guarantee}
 \begin{gathered}
 \frac{\zB}{\zK}\ge\frac{\zA}{\zK}\ge
 \rho_{\mathrm{par}}:=\sup_{t>0}\min_j\alpha_j(t)\ge f(D),\\
 f(D)=\begin{cases}
 \displaystyle\frac{2}{\sqrt{(1+D^2)^2+4D^2}+1+D^2},&0\le D\le1,\\
 ((1+\sqrt2)D)^{-1},&D\ge1.
 \end{cases}
 \end{gathered}
\end{equation}
The two branches agree at $D=1$. For $D\le1$, the simpler estimate
$f(D)\ge(1+2D^2)^{-1}$ also holds. In particular
$f(D)\ge\min\{1/3,0.41/D\}$ for $D>0$, and
$\zA=\zB=\zK$ as suprema when $D=0$.
\end{theorem}

\begin{proof}
Suppose first that $D>0$, and choose $t=\sqrt{\widetilde Y/\widetilde X}$.
Then $|b_j(t)|\le2D$. Write
\[
 g_j(s)=\frac{q(\sbar+s\widetilde p_j)}{q_0}
 =1+\ell_js+d_js^2,\qquad
 d_j=-\frac{\widetilde p_{jx}\widetilde p_{jy}}{q_0},\quad |d_j|\le D^2.
\]
Since $g_j(1)\ge0$, we have $\ell_j\ge-1-D^2$. If $D\ge1$, evaluating
at $1/D\le1$ instead gives $\ell_j\ge-D-d_j/D\ge-2D$.
For $\ell_j\ge0$, the denominator in \cref{dp:cylinder-step} is at most
$|b_j|\le2D$. For $\ell_j<0$ and $D\le1$, it is at most
\[
 \sqrt{(1+D^2)^2+4D^2}+1+D^2\le2+4D^2;
\]
the intermediate square-root bound follows by squaring, with difference
$8D^4$. For $D\ge1$, the denominator is at most $2D(1+\sqrt2)$.
These bounds prove the result for positive depth.

If $D=0$, all $d_j$ vanish. Feasibility and the scaled corner optimum give
$\min_j \ell_j=-1$: minimizing $\sum_j\lambda_j$ subject to
$1+\sum_j \ell_j\lambda_j\le0$ has value one. Let $t\to0$ if
$\widetilde Y=0$, or $t\to\infty$ if $\widetilde X=0$. Then every $b_j(t)$
tends to zero. Negative $\ell_j$ give limiting steps $-1/\ell_j$, while
nonnegative $\ell_j$ give infinite limiting steps. The finite-ray minimum tends
to one. Both family values are bounded above by $\zK$, proving their
supremum equality. This limiting argument does not assert attainment.
\end{proof}

In fixed coordinates, a Euclidean depth condition can imply this bound:
$\dist(\sbar,S)\le q_0$ by vertical projection, and
$\widetilde X\widetilde Y\le\operatorname{diam}(T^*)^2$. Consequently
$\delta\ge\dist(\sbar,S)/\operatorname{diam}(T^*)^2$ when the product is
positive. No condition on the ray matrix is needed in the theorem.

\begin{corollary}[A discriminant margin]\label{dp:grazing}
Suppose $D\ge1$. Assume that every ray missing $S$ on which $q$ initially
decreases has relative discriminant at most $-\mu$, where $0<\mu\le1$.
Put $\gamma=\sqrt{(1-\mu)/(1+\mu)}$ and
$\gamma_D=\max\{\gamma,1/D\}$. Then
\[
 \frac{\zA}{\zK}\ge\rho_{\mathrm{par}}\ge
 \frac{1}{D\bigl(\gamma_D+\sqrt{1+\gamma_D^2}\bigr)}.
\]
As $D\to\infty$, its leading constant ranges from $\sqrt2-1$ as
$\mu\downarrow0$ to one when $\mu=1$.
\end{corollary}

\begin{proof}
A decreasing restriction missing $S$ has $d_j>0$ and negative
discriminant. The margin gives $\ell_j^2\le4\gamma^2d_j$, hence
$-\ell_j\le2\gamma D$. A restriction meeting $S$ has its first positive
root at least one. If $d_j>0$, both roots are at least one and
$-\ell_j\le2$; if $d_j\le0$, the positive first root gives $-\ell_j\le1$.
Thus every negative $\ell_j$ satisfies $-\ell_j\le2D\gamma_D$.
Use $|b_j|\le2D$ in \cref{dp:cylinder-step}. Nonnegative $\ell_j$ give
at least $1/D$, which is at least the stated bound.
\end{proof}

\subsection{Sharp order, including completed sets}

The next family keeps its ray matrix fixed and makes the vertex approach the
surface. Its two horizontal rays miss $S$ with the largest possible negative
relative discriminant. These strong ray margins do not prevent loss.

\begin{proposition}[Fixed rays and vanishing depth]\label{dp:vanishing}
For $0<\varepsilon<1$, let
\[
 \sbar=(0,0,\varepsilon),\quad
 p_1=(1,-1,0),\quad p_2=(1,-2,0),\quad p_3=(1,1,1),\quad\rc=(1,1,1).
\]
Then $\zK=z_0=(1+\sqrt{1+4\varepsilon})/2$, uniquely attained on ray
three, with a transversal crossing. The relative discriminants are
$-1,-1,+1$, the ray matrix is fixed and nonsingular, and
$D=z_0\sqrt{2/\varepsilon}$. Nevertheless $\zB/\zK\to0$ as
$\varepsilon\downarrow0$. The lower bound in \cref{dp:depth-bound} applies.
\end{proposition}

\begin{proof}
Expansion gives
\[
 q(\sbar+P\lambda)=\varepsilon+\lambda_3-\lambda_3^2+
 \lambda_2\lambda_3+(\lambda_1+2\lambda_2)(\lambda_1+\lambda_2).
\]
Feasibility forces $\lambda_3\ge z_0$, and objective value $z_0$ forces
the other coordinates to vanish. The ray restrictions give the discriminants
and depth directly. The completed-set limit argument proving vanishing is
given in \cref{dp:vanishing-proof}. It is analytic and supplies no rate by
itself.
\end{proof}

\begin{theorem}[Exact computer-assisted sharpness bounds]\label{dp:certified-sharpness}
On the family in \cref{dp:vanishing}, for every $0<\varepsilon<1$,
\[
 \frac{\zA}{\zK}\le\frac{\zB}{\zK}\le
 \frac{137\sqrt\varepsilon}{z_0}=\frac{137\sqrt2}{D}.
\]
For
\[
 0<\varepsilon\le
 \varepsilon_1=\frac{53661932375105209}{2934348356061848092900},
\]
an explicit completed set, defined by a rational matrix after normalization,
gives
\[
 \zB/\zK\ge(1367/10)\sqrt\varepsilon/z_0.
\]

For $0<\eta<1$, $k\ge1$, and $L>0$, consider instead
\begin{equation}\label{dp:tangent-family}
 \sbar=(0,0,1),\quad p_1=(1,-1,-2(1-\eta)),\quad
 p_2=(1,-k,-2\sqrt k(1-\eta)),\quad p_3=(1,1,L),\quad \rc=(1,1,1).
\end{equation}
Here $L<\zK<L+1/L$ and $D=\sqrt k\,\zK$. The following exact bounds hold;
the upper bounds are computer-assisted, and the lower bounds are obtained
from printed rational witnesses:
\begin{center}
\begin{tabular}{@{}cccc@{}}\toprule
 $\eta$ & $k$ & $D\zB/\zK$ upper, all $L>0$ &
 $D\zB/\zK$ lower, $L\ge17/5$\\\midrule
 $1/1000$ & $4$ & $197/100$ & $2461/1250$\\
 $1/100$ & $4$ & $5/2$ & $6117/2500$\\
 $1/1000$ & $9/4$ & $63/40$ & $15621/10000$\\
 $1/1000$ & $49/25$ & $77/50$ & $18851/12500$\\\bottomrule
\end{tabular}
\end{center}
Every upper bound also bounds $D\zA/\zK$.
\end{theorem}

\begin{proof}
For \cref{dp:tangent-family}, put $m=\lambda_1+\sqrt k\lambda_2$.
Cauchy--Schwarz gives
$(\lambda_1+\lambda_2)(\lambda_1+k\lambda_2)\ge m^2$, and expansion yields
\[
 q\ge (m-(1-\eta))^2+2\eta-\eta^2+L\lambda_3-\lambda_3^2.
\]
Thus feasibility forces $\lambda_3>L$, while ray three alone reaches $S$
before $L+1/L$. The two missing ray restrictions have minimum
$2\eta-\eta^2$ and relative discriminant
$((1-\eta)^2-1)/((1-\eta)^2+1)$. The third discriminant is $+1$.
The upper bounds and rational lower witnesses are established in
\cref{dp:box-proof,dp:lower-witnesses}.
\end{proof}

For each family separately, define its asymptotic worst-case constant by
\[
 C_{\famA}^*=\lim_{D_0\to\infty}\inf_{D\ge D_0}
             D\frac{\zA}{\zK},\qquad
 C_{\famB}^*=\lim_{D_0\to\infty}\inf_{D\ge D_0}
             D\frac{\zB}{\zK},
\]
where the infima range over the corners in this section. The limits exist
by monotonicity. The analytic lower bound and the exact computer-assisted
upper bound imply
\[
 \sqrt2-1\le C_{\famA}^*\le C_{\famB}^*\le\frac{77}{50}=1.54.
\]
Thus the order $1/D$ is sharp for both families, and the analytic constant
is within a factor $1.54(1+\sqrt2)<3.72$ of either worst-case constant.
The constants need not be equal. Bounded depth gives a uniform guarantee;
large depth permits bad corners, while some large-depth corners remain exact.

\subsection{A fixed point rule and its coordinate dependence}

Consider the Case-4 rule applied to the constraint exactly as
$w-xy\le0$: unit coefficient on $w$, no additional linear terms in $x,y$,
and source constant $\kappa_{\mathrm{SCIP}}=0$. All coordinates below are
those of this representation. The rule chooses the unit vector proportional
to $((\bar x-\bar y)/2,(\bar w+1)/2)$ in the positive Sylvester coordinates.
Set
\[
 N_0^2=(\bar x-\bar y)^2+(\bar w+1)^2,\qquad
 V(d)=(\bar x-\bar y)d_w-(\bar w+1)(d_x-d_y).
\]
The uncompleted orbit set is
\begin{equation}\label{dp:fixed-rule-set}
 C_{\mathrm{fix}}=
 \left\{s:4N_0^2q(s)\ge V(s-\sbar)^2,
 (\bar w+1)(w+1)+(\bar x-\bar y)(x-y)\ge0\right\}.
\end{equation}
The actual completed set is $B_{\mathrm{fix}}=C_{\mathrm{fix}}+\R_+e_w$.
Write $z_{\mathrm{fix}}$ for its cut value. The first positive root of
$4N_0^2q(\sbar+sp)=s^2V(p)^2$ gives the uncompleted step. A completed
step must instead allow vertical lowering; it can be larger or infinite.
The rotation derivation and an explicit example distinguishing the two
steps are in \cref{dp:fixed-rule-derivation}.

\begin{theorem}[Fixed-rule guarantee]\label{dp:fixed-rule-bound}
For any finite ray family, both the uncompleted and completed rule values
satisfy
\[
 \frac{z_{\mathrm{fix}}}{\zK}\ge
 \min\left\{\frac12,\frac1{2D},
 \frac{\sqrt{q_0}}{\max_j\norm{(\widetilde p_{jx}-\widetilde p_{jy},
                                  \widetilde p_{jw})}}\right\}.
\]
Reciprocals of zero are interpreted as infinite. If the $3\times N$ scaled
ray matrix has full row rank and spectral condition number $\kappa$, then
\[
 \frac{z_{\mathrm{fix}}}{\zK}\ge
 \min\left\{\frac12,\frac1{\sqrt{2N}\,\kappa D}\right\}.
\]
In particular the constant is $1/\sqrt6$ for three independent rays.
\end{theorem}

\begin{proof}
Put $s_0=\min\{1/2,1/(2D)\}$. The polynomial $g_j$ used above is at least
$1/4$ on $[0,s_0]$. Indeed, for $d_j\le0$, concavity gives
$g_j(s)\ge1-s$; for $d_j>0$ with a forward root, both roots are at least
one and $g_j(s)\ge(1-s)^2$. If there is no forward root and $\ell_j<0$,
$\ell_j>-2\sqrt{d_j}$, so $g_j(s)\ge(1-\sqrt{d_j}s)^2\ge(1-Ds)^2$
on this interval. If $\ell_j\ge0$, the assertion is immediate.
The determinant inequality in \cref{dp:fixed-rule-set} therefore holds
through the additional length $N_0\sqrt{q_0}/|V(\widetilde p_j)|$.
Since
$|V(p)|\le N_0\norm{(p_x-p_y,p_w)}$, this proves the first bound.
The PSD matrix starts positive definite at $\sbar$ and stays so before its
first determinant zero; hence the trace branch cannot change first.

For the second bound,
$\norm{(p_x-p_y,p_w)}\le\sqrt2\norm p\le\sqrt2\norm{\widetilde P}_2$.
Full row rank gives
$\sigma_{\min}(\widetilde P)\le\norm{\operatorname{row}_x(\widetilde P)}
\le\sqrt N\widetilde X$, and likewise with $y$. Hence
$\sigma_{\min}\le\sqrt N D\sqrt{q_0}$ and
$\norm{\widetilde P}_2\le\kappa\sqrt N D\sqrt{q_0}$.
Since full row rank requires $N\ge3$ and $\kappa\ge1$, the term
$1/(\sqrt{2N}\kappa D)$ also does not exceed $1/(2D)$.
Upward completion can only increase the cut value.
\end{proof}

\begin{proposition}[Sharp loss in $\kappa D$]\label{dp:fixed-rule-sharp}
Given any $\kappa\ge1$ and $D>0$, put $v=\kappa D$,
$u=\sqrt{1+v^2}$, and take unit costs and
\[
 \sbar=(u,-u,-1),\quad
 p_1=(vD,0,0),\quad p_2=(0,-vD,0),\quad p_3=(0,0,-v^2).
\]
This corner has exactly depth $D$, condition number $\kappa$, and
$\zK=\zA=\zB=1$, with the family optimum attained. All three relative
discriminants are $+1$. Nevertheless both fixed-rule sets give
\[
 \frac{z_{\mathrm{fix}}}{\zK}=
 \frac{2}{1+\sqrt{1+\kappa^2D^2}}<\frac2{\kappa D}.
\]
Thus for three rays the worst fixed-rule ratio at prescribed $\kappa,D$
has order $\min\{1,1/(\kappa D)\}$.
\end{proposition}

The complete calculation is in \cref{dp:fixed-rule-sharp-proof}.
A simpler rescaling obstruction, proved in \cref{dp:rescaling}, keeps
$D=1$ and an exact cylinder while the fixed-rule value tends to zero.
The constant $+1$ in the point rule introduces a coordinate scale. Changing
the constraint representation, for example to $2w-2xy\le0$ or
$w-xy-c\le0$, changes the rule. These statements therefore do not apply
unchanged after rescaling, shifting, or to the two-variable Case-2 rule.
They concern the set geometry before numerical safeguards or additional
strengthening, and do not imply a solver performance improvement.

\subsection{Contact geometry: supremum and attainment}

Suppose now that the corner optimum has a unique contact point $t^*$ with
$q(t^*)=0$. Translate the bilinear coordinates by a symmetry so that
$t^*=0$. For a vertex $v$ of $T^*$, its new coordinates are
$(u_{vx},u_{vy},h_v)$, where
\[
 u_v=v-t^*,\qquad h_v=\nabla q(t^*)^Tu_v,
 \qquad q(v)=h_v-u_{vx}u_{vy}.
\]
The height $h_v$ measures transversality to the tangent plane. Assume that
$t^*$ is one of the scaled-ray vertices and $h_v>0$ for every other vertex.
Uniqueness then implies $q(v)>0$ at those other vertices.

Every orbit set through the contact point has the form
\begin{equation}\label{dp:contact-set}
 C_{\alpha,\beta}=
 \{(x,y,w):4\alpha q\ge(\alpha x-y-\beta w)^2,
                 \alpha w+\beta x+1\ge0\},\qquad \alpha>0, \beta\in\R.
\end{equation}
A completion containing the contact has a generating orbit set through it:
the allowable lowering at a point with $q=0$ is zero.
For a point with positive height $h$, membership is equivalent to
$\beta\in I_v(\alpha)$, where
\begin{equation}\label{dp:interval}
 I_v(\alpha)=
 \left[\frac{\alpha u_{vx}-u_{vy}-2\sqrt{\alpha q(v)}}{h_v},
       \frac{\alpha u_{vx}-u_{vy}+2\sqrt{\alpha q(v)}}{h_v}\right].
\end{equation}
Thus all noncontact vertices impose intervals on a common scalar parameter.

\begin{lemma}[Approximation from a closed contact set]\label{dp:contact-approx}
If an orbit set contains $T^*$, with $\sbar$ possibly on its boundary, then
$\zA=\zK$ as a supremum. No rank assumption is needed for this implication.
\end{lemma}

\begin{proof}
Write its matrix as $G$ with $\det G>0$, and fix $0<r<1$. Set
$G_\epsilon=G+\epsilon M(\sbar)^{-1}$. At $\sbar$ this adds $\epsilon I$
to the PSD matrix. At any vertex $v_r=(1-r)\sbar+rv$ of $T_r$, a null
vector of its PSD matrix is a common null vector of the endpoint matrices.
Along the segment, write $GM(s)=\sym(GM(s))+k(s)J$, with
$J=\left(\begin{smallmatrix}0&1\\-1&0\end{smallmatrix}\right)$.
On such a nonzero common null vector $z$,
$GM(s)z=k(s)Jz$. Since $q(s)>0$ on the contracted segment, $M(s)$ is
invertible and $k(s)$ cannot vanish. It retains its sign. Therefore
\[
 M(\sbar)^{-1}M(v_r)z=\frac{k(v_r)}{k(\sbar)}z,
\]
so the perturbation is positive on the PSD kernel. This makes every
contracted vertex matrix positive definite for one sufficiently small
positive $\epsilon$, while $\det G_\epsilon>0$ by continuity.
Hence an admissible orbit set contains $T_r$ and has cut ratio at least $r$.
Let $r\uparrow1$. Full details of the kernel and perturbation step are
given in \cref{dp:contact-proof}.
\end{proof}

\begin{theorem}[Closed and strict contact intervals]\label{dp:contact-intervals}
Under the unique support-one contact and positive-height assumptions above,
\begin{enumerate}
 \item $\zA=\zK$ as a supremum if and only if, for some $\alpha>0$,
       all intervals $I_v(\alpha)$, $v\ne t^*$, have a common point;
 \item a member of $\famA$ attains $\zK$ if and only if such a common point
       lies in $\intt I_{\sbar}(\alpha)$;
 \item if no closed intersection exists for any $\alpha>0$, then
       $\zA<\zK$.
\end{enumerate}
\end{theorem}

\begin{proof}
An orbit set through zero has
$G=\left(\begin{smallmatrix}\alpha&0\\\beta&1\end{smallmatrix}\right)$
after positive scaling, by rank-one PSD membership at $M(0)=e_2e_2^T$.
Its determinant and trace give \cref{dp:contact-set}.
For $h>0$, the first diagonal entry is $\alpha h>0$, so the determinant
inequality already implies the trace branch. Solving it for $\beta$ gives
\cref{dp:interval}. Strict membership of the vertex is precisely interior
membership in its interval. This proves the attainment criterion, and
\cref{dp:contact-approx} proves the sufficiency for supremum equality.

For necessity, normalize an exactness-approximating sequence of matrices
and take a nonzero limit $G$. Its determinant is nonnegative, and its PSD
inequalities contain $T^*$. If the determinant is positive, the intervals
follow as above. If it is zero, PSD membership at zero gives
$G=\left(\begin{smallmatrix}a&0\\c&d\end{smallmatrix}\right)$ with
$d\ge0$ and $ad=0$. If $d>0$, then $a=0$ and PSD membership at every
positive-height vertex forces $y_v=-(c/d)h_v$. Set $\beta=c/d$ and
choose $\alpha>0$ sufficiently small. The determinant slack is
$4\alpha q(v)-\alpha^2x_v^2>0$ at every noncontact vertex.
If $d=0$, PSD membership at $\sbar$ forces $a>0$, and the determinant
equation forces $x_v=(c/a)h_v$. Set $\beta=\alpha c/a$ and choose
$\alpha$ sufficiently large. The slack is then
$4\alpha q(v)-y_v^2>0$. Uniqueness gives $q(v)>0$, and finiteness gives
one common choice of $\alpha$. Either singular case therefore produces
an attaining orbit set. This proves necessity and the strict-gap assertion.
\end{proof}

The use of closed intervals in the first part is essential. They characterize
an exact supremum, while strict membership at the LP vertex characterizes an
attaining set. The theorem concerns $\famA$; allowing independent vertical
lowering in $\famB$ requires further analysis.

\begin{corollary}[A quantitative cylinder condition]\label{dp:cylinder-contact}
If, for some $t>0$,
\[
 4h_v\ge(tu_{vx}+u_{vy}/t)^2\quad(v\ne t^*),
\]
with strict inequality at $\sbar$, then a parabolic cylinder through $t^*$
attains $\zA=\zB=\zK$. In particular this holds when
\[
 H_{\mathrm{cyl}}:=\sup_{t>0}\min_{v\ne t^*}
   \frac{4h_v}{(tu_{vx}+u_{vy}/t)^2}>1,
\]
where a zero denominator means an infinite ratio. A cruder sufficient
condition is $h_v\ge XY$ at every other vertex, strictly at $\sbar$,
where $X=\max_{v\ne t^*}|u_{vx}|$ and $Y=\max_{v\ne t^*}|u_{vy}|$.
\end{corollary}

\begin{proof}
Set $\beta=0$ and $\alpha=t^2$ in \cref{dp:contact-set}.
Its determinant inequality is exactly
$(tu_x+u_y/t)^2\le4h$. The set is an upward-closed cylinder, so it
belongs to both families. For $X,Y>0$, take $t=\sqrt{Y/X}$ and use
$(t|u_x|+|u_y|/t)^2\le4XY$. If $XY=0$, finite positive heights allow
a sufficiently small or large finite $t$ instead.
\end{proof}

\begin{proposition}[Sharp cylinder threshold]\label{dp:cylinder-sharp}
For every $0<d<7/128$, take unit costs, $t^*=0$, and scaled-ray vertices
\[
 \sbar=(3/2,-1/2,1/4-d),\quad
 v_2=(1/2,3/2,1-d),\quad v_3=(3,1,4-d),
\]
with rays from $\sbar$ to $t^*,v_2,v_3$. Then $\zK=1$, the unique
minimizer is $t^*$, all incident edges have positive tangent height, and
$H_{\mathrm{cyl}}\ge1-4d$, but $\zA<1$.
Thus no threshold below one can replace the strict sufficient threshold in
\cref{dp:cylinder-contact}.
\end{proposition}

\begin{proof}
At $d=0$, the minimum of $q$ on the opposite triangular face is $7/128$,
and its height is at least $1/4$; the explicit face calculation is in
\cref{dp:cylinder-sharp-proof}. For a point $\theta v$ toward zero,
\[
 q_d(\theta v)/\theta=(1-\theta)w_0(v)+\theta q_0(v)-d
 \ge7/128-d>0\qquad(0<\theta\le1).
\]
Hence zero is the unique feasible point of $T^*$ and $\zK=1$.
At $t=1$, the three height ratios are $1-4d,1-d,1-d/4$.
The explicit interval separation in \cref{dp:cylinder-sharp-proof}
holds for all $0<d\le1/8$ and rules out every closed intersection.
Apply \cref{dp:contact-intervals}.
\end{proof}

\begin{proposition}[Angles alone do not ensure exactness]\label{dp:angles}
For every $\theta<\pi/2$, there is a support-one corner with $\zA<\zK$ whose
edges at its unique contact point all make angles at least $\theta$ with
the tangent plane.
\end{proposition}

\begin{proof}
Use any instance of \cref{dp:cylinder-sharp}, centered at its contact.
The symmetry $(x,y,w)\mapsto(kx,ky,k^2w)$ preserves the family values,
depth, and $H_{\mathrm{cyl}}$. An edge $(u_x,u_y,h)$ with $h>0$ has sine
of its angle with the tangent plane equal to
\[
 \frac{k^2h}{\sqrt{k^2(u_x^2+u_y^2)+k^4h^2}},
\]
which tends to one. There are finitely many edges, so one $k$ suffices.
\end{proof}

For a two-sided tangent edge through $t^*$ with direction $d$ satisfying
$d_w=0$ and $d_xd_y<0$, PSD membership forces
$\alpha=-d_y/d_x$ in \cref{dp:contact-set}. This is the tangent pencil
in \cref{fd:tangent-pencil}. Positive-height vertices then impose the
same intervals with $\alpha$ fixed. Zero-height edge points must also
satisfy $\alpha x+y=0$ and $\beta x+1\ge0$; they cannot be checked by
dividing by their height. This distinguishes tangent rigidity from the
strictly positive-height support-one criterion above.
